\begin{abstract}

\vspace{0.6em}
\noindent Anomaly detection in surveillance camera videos is a cornerstone of public safety, yet real deployments face three obstacles. First, general vision-language models suffer a severe domain shift from the surveillance domain, which systematically depresses anomaly scores and prevents thresholds from transferring across scenes. Second, explanations are generated freely by a large model, and there is no way to tell which claims are actually supported by video evidence. Third, training-free pipelines implicitly assume clean inputs: low light, rain and fog, video compression and camera shake collapse their scores and induce hallucinated semantics. We present \evivad{}, a unified framework built on a \emph{fully frozen} vision-language backbone with three complementary modules. (1) \textbf{DAA} (Domain Low-rank Adapter) injects rank-$r{=}4$ low-rank increments into the Q/V projections of the last four encoder blocks, adapting to the surveillance domain with fewer than $5$\,M trainable parameters; the up-projection is zero-initialised so that the starting point is bit-wise identical to the baseline. (2) \textbf{EAD} (Evidence-Anchored Decoding) forces explanations to explicitly cite temporal intervals and visual cues, and is trained with an InfoNCE evidence-alignment loss, so that an explanation can be both supported by evidence and falsified by counterexamples; when evidence is insufficient the model refuses to answer. (3) \textbf{DAG} (Degradation-Aware Gating) uses a lightweight quality probe to blend a vision-text alignment path with a retrieval/prior-driven path, and abstains under low confidence instead of emitting confident false alarms. We establish a unified protocol on UCF-Crime, three cross-domain sets (XD-Violence, UBnormal, MSAD) and a $4\times3$ degradation grid. The aggregate results indicate that, relative to the training-free baseline B0, \evivad{} raises frame-level AUC from $76.8$ to $82.1$, improves evidence-attribution recall (EAR) from $41.5$ to $68.4$, reduces the hallucination rate (HR) from $26.4$ to $9.8$, and lifts the average AUC on the degradation grid from $62.5$ to $71.3$ (relative performance retention $0.81 \rightarrow 0.92$), at a cost of only $5.5$\,M additional trainable parameters. Our main message is that in the era of large models, video anomaly detection should not be evaluated by detection scores alone: the \emph{verifiability} of explanations and the \emph{graceful degradability} under input corruption deserve equal standing.
\end{abstract}
